\chapter*{Roadmap of the Theory}\addcontentsline{toc}{chapter}{Roadmap of the Theory}
\markboth{ROADMAP OF THE THEORY}{}

The theory of this book follows a single chain of reasoning, from
the concrete problem of encryption to the abstract geometry of the
Universal Optimality Defect.  The diagram below shows the complete
logical flow.

\vspace{1em}

\begin{center}
\begin{tikzpicture}[
  node distance=1.2cm,
  arrow/.style={->,>=stealth,thick},
  box/.style={rectangle,draw,rounded corners=2pt,
    minimum width=8cm,minimum height=0.7cm,
    align=center,font=\small}
]

\node[box] (enc) {{Universally Optimal Encoder}};
\node[box,below=of enc] (mult) {Multiplicity Matrix $n_{mc}$};
\node[box,below=of mult] (wis) {Weighted Incidence Structure $(w,\beta)$};
\node[box,below=of wis] (gauge) {Gauge Class $[w,\beta]$};
\node[box,below=of gauge] (phi) {Mismatch Operator $\Phi=[A\;-C]$};
\node[box,below=of phi] (lsq) {{Least-Squares Problem $\min\|Au-\omega\|^2$}};
\node[box,below=of lsq] (norm) {Normal Operator $\Phi^{\mathsf T}\Phi$};
\node[box,below=of norm] (lap) {Posterior Laplacian $L=A^{\mathsf T}QA$};
\node[box,below=of lap] (quot) {Quotient Geometry $V_I/\operatorname{Im}(C)$};
\node[box,below=of quot,fill=gray!10] (uod) {{Universal Optimality Defect $\mathcal D(P)$}};

\draw[arrow] (enc) -- node[right] {\small Chapter~\ref{ch:sec:background}} (mult);
\draw[arrow] (mult) -- node[right] {\small Chapter~\ref{ch:sec:wis}} (wis);
\draw[arrow] (wis) -- node[right] {\small Chapter~\ref{ch:sec:representation}} (gauge);
\draw[arrow] (gauge) -- node[right] {\small Chapter~\ref{ch:sec:graphs}} (phi);
\draw[arrow] (phi) -- node[right] {\small Chapter~\ref{ch:sec:hilbert-geometry}} (lsq);
\draw[arrow] (lsq) -- node[right] {\small Chapter~\ref{ch:sec:hilbert-geometry}} (norm);
\draw[arrow] (norm) -- node[right] {\small Chapter~\ref{ch:sec:the-posterior-manifold}} (lap);
\draw[arrow] (lap) -- node[right] {\small Chapter~\ref{ch:sec:canonical-projection-and-differential-factorisation}--\ref{ch:sec:local-diffeomorphism}} (quot);
\draw[arrow] (quot) -- node[right] {\small Chapter~\ref{ch:sec:pullback-riemannian-metric}} (uod);

\end{tikzpicture}
\end{center}

\vspace{1em}

\noindent\textbf{How to read this diagram.}
Start at the top: a universally optimal encoder produces a matrix of
multiplicities $n_{mc}$ (Chapter~\ref{ch:sec:background}).  This matrix factorises into a
weighted incidence structure (Chapter~\ref{ch:sec:wis}), whose only ambiguity is a
global scaling---the gauge class (Chapter~\ref{ch:sec:representation}).  Taking logarithms
linearises the structure, producing the mismatch operator $\Phi$
(Chapter~\ref{ch:sec:graphs}), which leads to a least-squares problem whose solution
involves the normal operator $\Phi^{\mathsf T}\Phi$ (Chapter~\ref{ch:sec:hilbert-geometry}).
The Schur complement~\cite{zhang2005schur} of this operator is the posterior Laplacian $L$
(Chapter~\ref{ch:sec:the-posterior-manifold}), which encodes the geometry of the quotient space
$V_I/\operatorname{Im}(C)$ (Chapters~\ref{ch:sec:canonical-projection-and-differential-factorisation}--\ref{ch:sec:local-diffeomorphism}).  The distance from a
distribution to the optimal set---the Universal Optimality Defect
$\mathcal D(P)$---is the squared norm in this quotient geometry
(Chapter~\ref{ch:sec:pullback-riemannian-metric}).

\vspace{1em}

\noindent Each subsequent part of the book builds on this core chain:
Part~IV (Chapters~\ref{ch:sec:entropy-bounds}--\ref{ch:sec:universal-security-bounds}) uses this geometry to bound information
functionals; Part~V (Chapters~\ref{ch:sec:analytic-uod}--\ref{ch:sec:open}) develops the variational
structure of the defect and its comparison with existing frameworks.
